\section{Attainment on the original smooth type}\label{sec:attainment}

\begin{theorem}[Fixed-type attainment, \cite{Baker2026BourAttainment}]
\label{thm:attainment}
If \(I=I^+(M)<2T\), there is a smooth positive Bach-flat class on
the oriented \(M\) with \(E_+=E_-=I\). On a nonstandard entered
type, \(I=a_M\), the least energy of a positive Bach-flat class in
the strict subcritical band.
\end{theorem}

\begin{proof}[Argument of the companion theorem]
Fix the smooth type first. Choose \(0<\alpha_j\downarrow0\) and
volume-one positive Yamabe metrics on \(M\) with initial excess
\(\epsilon_j=e_{j,0}-I\in[0,\alpha_j^2]\). For fixed \(j\) suppress
the subscript. Bour's functional in our norms is
\[
 {\cal F}^\alpha=2(1-\alpha)e+\tfrac{\alpha}{2}D
                =2e+\tfrac{\alpha}{24}\int R^2\,dV-16\alpha\pi^2 .
\]
The initial value is at most \(2e_0\) by \eqref{eq:budgets}. For late
\(j\), \(e_0/(1-\alpha)\le2T-\eta_0\) for some fixed \(\eta_0>0\).
In particular Bour's pinching hypothesis is strict. Along its negative gradient flow monotonicity gives
\(e(t)\le e_0/(1-\alpha)\) at every smooth time on the original flow.
By \eqref{eq:zero-wall} and \eqref{eq:budgets}, each slice is in a
positive class with \(Y\ge\sqrt{48\eta_0}\), and hence \(e(t)\ge I\).
On the physical flow metric one therefore has
\begin{equation}\label{eq:original-flow}
 D(t)\le D_0+\frac{4(1-\alpha)}{\alpha}(e_0-e(t))
       \le 4I+\frac{4\epsilon_j}{\alpha_j}
       \le 4I+4\alpha_j<32\pi^2 .
\end{equation}
The last inequality has a fixed positive margin for late \(j\).
Here \(D_0\le4e_0\) was used in the starting Yamabe gauge. The second
formula for \({\cal F}^\alpha\) also gives
\(\int R(t)^2\,dV_{g(t)}\le Y_0^2+48\epsilon_j/\alpha_j\), where
\(Y_0\) is the initial Yamabe constant.

\emph{The two outcomes at fixed \(\alpha\).}
Bour's Theorem~1.11 gives a closed \({\cal F}^\alpha\)-critical
subsequential limit or a complete noncompact nonflat scalar-flat
Bach-flat limit at a curvature maximum
\cite[Theorems 1.4, 1.9 and 1.11]{Bour2010FourthOrderCurvatureFlows}.
The companion justifies the short-time uniqueness and bounded-slab
smoothing used in this strict case \cite{BahuaudHelliwell2015Uniqueness,Baker2026BourAttainment}.
The closed endpoint stays on \(M\) and has constant scalar curvature
by the trace of the regularized equation. It can have nonzero Bach tensor.

For a pointed limit \(X\), compactly supported tests and compact
exhaustion pass the Yamabe lower bound, each chiral upper bound and
\eqref{eq:original-flow} to the limit. The homogeneous Sobolev estimate
and finite curvature give finitely many asymptotically locally Euclidean
ends with differentiated decay
\cite{TianViaclovsky2005ALE,TianViaclovsky2008Volume,AcheViaclovsky2011Obstruction}.
On a connected neighborhood of the \emph{whole} truncated core, including
the buffers of every end, take one oriented comparison embedding into the
same source slice. Its disjoint ends can all be straightened to round
quotient products on the still closed source with arbitrarily small
chiral cost. The spherical-link obstruction at \(E_+<9\pi^2\) forces
each quotient group to be trivial \cite[A quotient obstruction on the
closed source]{Baker2026MarkedNecks}. If there are \(k\) ends, capping
the core gives a homotopy sphere by the marked-cut topology, so
\(\chi(X)=2-k\). The Euler end term can be computed using only the closed formula.
On each dyadic end annulus splice the differentiated ALE metric to the
Euclidean metric. In the radius-one rescaling the change of its total
quadratic Gauss--Bonnet integral is \(o(1)\). Close each resulting flat
outlet of radius \(r\) with a fixed scaled ball whose smooth warped
metric continues the increasing flat radius across the seam and closes
at a round pole. This cap contributes exactly \(16\pi^2\): joined to the
flat radius-\(r\) ball it gives a smooth \(S^4\), and the flat ball
contributes zero. The capped core has Euler characteristic two.
Letting \(r\) tend to infinity in its closed formula gives
\[
 D_X=2\big(E_+(X)+E_-(X)\big)+32\pi^2(k-1).
\]
The strict bound in \eqref{eq:original-flow} forces \(k=1\).

Conformal compactification and removal of its point give a smooth
positive Bach-flat homotopy sphere \(N\). In particular its chiral
energies agree and \(D_X=4e(N)\). The closed critical gap and
nonflatness imply \(e(N)>T\)
\cite[Strict recognition and the closed critical gap]{Baker2026MarkedNecks}.
The smooth type still has to be identified. Remove a small normal ball
in the smooth compactifying chart and embed its complement with both
collar buffers into one original slice. In each sign the upper limit of the whole raw
exterior energy is at most \(2T-\eta_0-e(N)+o(1)<T\).
Choose the buffer length first, then the ball radius, then the comparison
index. The buffered cap estimate gives a positive cap on that actual
exterior of energy below \(T\). Recognition makes this cap a standard
sphere. Disk-neighborhood extension identifies the original exterior
with the deleted normal ball, matching the inverse core map on their
common collar \cite[Theorem C]{Palais1960Extending}. Gluing the two
maps gives an oriented diffeomorphism \(M\cong N\), with smooth
inverse in the same marking. In particular \(I\le e(N)\) and exhaustion
gives \(e(N)\le e_0/(1-\alpha)\).
Keeping the initial Yamabe term in \eqref{eq:original-flow} gives
\(4(e(N)-I)+(Y_{\rm rd}^2-Y_0^2)/12\le4\epsilon_j/\alpha_j\).
Thus pointed outcomes along infinitely many original flows require
their initial Yamabe constants to approach the round value.

\emph{The limit as \(\alpha_j\to0\).}
Retain the volume-one physical representative of a closed regularized
endpoint. For a compactified pointed endpoint take its volume-one
Yamabe representative on the identified \(M\) and recompute
\(D\le4e\) there. The resulting closed metrics \(q_j\) have
\(I\le e(q_j)\le e_{j,0}/(1-\alpha_j)\), uniformly positive Yamabe
constants, bounded constant scalar curvatures and a common strict
Ricci-energy bound below \(32\pi^2\). In the constant-scalar Euler
equation the leading term is \(\tfrac12\nabla^*\nabla Z\), and the
parameter \(p_j=\alpha_j\) or \(0\) enters only in a bounded zeroth-order
coefficient. On curvature-normalized balls the Yamabe inequality gives
noncollapse. In harmonic charts the Ricci equation for the metric and
the elliptic equation for \(Z\) give uniform interior derivative bounds
\cite[The limit of critical metrics]{Baker2026BourAttainment}. These are
stationary estimates uniform in \(j\).

Bounded curvature now gives a closed smooth Bach-flat limit on \(M\)
with energy \(I\). At unbounded curvature a rescaling about a maximum
gives another complete nonflat scalar-flat Bach-flat limit. The same
marked compactification argument applies to this sequence of \(q_j\).
Its positive class on \(M\) satisfies
\(I\le e(N)\le\liminf_j e(q_j)=I\). In this last case the exterior mass in either sign tends to zero when the source index tends to infinity first and the normal ball radius then tends to zero. For a pointed output at a fixed positive parameter the corresponding upper bound is \(e_0/(1-\alpha)-I\). This proves attainment in both
cases. On a nonstandard type the minimum lies above \(T\). Every positive
exact class is admissible for \(I\), and the attaining class is itself
exact, proving \(a_M=I\).
\end{proof}
